\section{La esencia de las redes neuronales: aproximadores de funciones}

\subsection{Universal Approximation Theorem}

Para entender el poder y las limitaciones del aprendizaje profundo, un punto de partida muy importante es el \textbf{Universal Approximation Theorem} (teorema de aproximación universal) propuesto por Hornik y colegas en 1989:

\begin{importantbox}{Universal Approximation Theorem}
Una red neuronal prealimentada de una sola capa oculta, si la cantidad de unidades ocultas es suficientemente grande, puede aproximar cualquier función continua con una precisión arbitraria.
\end{importantbox}

$$
f(x) \approx \hat{f}(x; W, b) = \sigma(W_2 \cdot \sigma(W_1 x + b_1) + b_2)
$$

Donde:
\begin{itemize}
    \item $f(x)$: función continua objetivo
    \item $\hat{f}(x; W, b)$: salida de la red neuronal con una sola capa oculta
    \item $\sigma$: función de activación no lineal
    \item $W_1, W_2, b_1, b_2$: parámetros de pesos y sesgos aprendibles
\end{itemize}

\begin{figure}[H]
\centering
\includegraphics[width=0.85\textwidth]{slides-images/slide-12.jpg}
\caption{Diagrama esquemático del Universal Approximation Theorem\protect\footnotemark}
\end{figure}
\footnotetext{Diapositiva 12. Cita a Hornik et al., Neural Networks, 1989.}

\subsection{Las limitaciones del teorema}

Aunque este teorema es muy potente en el plano teórico, existen dos caveats clave:

\begin{figure}[H]
\centering
\includegraphics[width=0.85\textwidth]{slides-images/slide-13.jpg}
\caption{Dos advertencias importantes del Universal Approximation Theorem\protect\footnotemark}
\end{figure}
\footnotetext{Diapositiva 13.}

\begin{warningbox}{Las dos advertencias principales del teorema de aproximación universal}
\begin{enumerate}
    \item \textbf{La cantidad de unidades ocultas puede ser inviablemente grande}: el teorema garantiza la existencia de dicha red, pero no da una cota superior para la cantidad de neuronas necesarias: para funciones complejas, podría requerirse una cantidad astronómica de parámetros.
    \item \textbf{No garantiza la capacidad de generalización}: el teorema solo garantiza la capacidad de aproximación, no que el modelo aprendido se desempeñe bien con datos no vistos. Ajustarse a los datos de entrenamiento no equivale a comprender las regularidades subyacentes de los datos.
\end{enumerate}
\end{warningbox}

Además, el teorema no ofrece una metodología para encontrar los pesos óptimos. Si bien el descenso de gradiente es una estrategia de optimización eficaz, en superficies de pérdida no convexas de alta dimensión, la optimización en sí misma es un problema sumamente no trivial.

\subsection{Perspectiva histórica del desarrollo de la IA}

\begin{figure}[H]
\centering
\includegraphics[width=0.85\textwidth]{slides-images/slide-14.jpg}
\caption{La trayectoria histórica del ``Hype'' de la IA\protect\footnotemark}
\end{figure}
\footnotetext{Diapositiva 14.}

El desarrollo de la IA atravesó un hype cycle típico: desde el nacimiento en la conferencia de Dartmouth en 1956, pasando por el primer invierno de la IA de 1974 a 1980, luego el segundo invierno de 1987 a 1993, hasta llegar finalmente a un crecimiento explosivo en la era del aprendizaje profundo. Comprender esta historia nos ayuda a mantener un juicio sereno sobre el auge actual de la IA: hay que abrazar el avance tecnológico, pero también estar alerta ante el exceso de exaltación.

\subsection{Resumen de la sección}

En esencia, las redes neuronales son aproximadores de funciones. El teorema de aproximación universal garantiza la capacidad expresiva en el plano teórico, pero en la aplicación práctica, el tamaño de la capa oculta, la dificultad de la optimización y la capacidad de generalización son desafíos reales que hay que enfrentar. Este marco de comprensión sienta las bases para la discusión, en la siguiente sección, de las limitaciones concretas del aprendizaje profundo.

%% =====================================================
